\section{A cycle-count recurrence without witness chords}
\label{sec:cycle-recursion}

Fable 5.2 proposed this use of the Entringer--Slater deletion identity in its November 10 notes.
The recurrence gives a shorter route from a one-cycle forest estimate to the cycle-count asymptotic theorem.
It does not change the existing formal proof.
It also removes the separate witness-avoidance error from the quantitative argument.

Let $\psi(r)$ be the maximum number of cycles among finite graphs of rank $r$.
The same maximum results if one requires simple connected graphs.
Bridges can join different components, and subdividing each edge twice removes loops and parallel edges without changing cycles or rank.
Since cycle counts are integers bounded by $2^r-1$, this maximum is attained.

\begin{lemma}[One marked edge]\label{lem:recursion-marked-edge}
Let $G$ be a connected multigraph of rank $r$, let $W$ be a cycle, and let $e\in E(W)$.
Then
\[
 c_e(G)\le p(G,W)2^{r-1}.
\]
\end{lemma}

\begin{proof}
A cycle $C$ containing $e$ has a forest outside $W$.
If two such cycles $C,C'$ have the same outside part, their symmetric difference is an even edge set in $W-e$.
Since $W-e$ is a forest, that symmetric difference is empty.
Thus these cycles give distinct forest outside parts.
There are $2^{r-1}$ outside parts, since $W$ has rank one.
Their forest proportion is $p(G,W)$, proving the bound.
\end{proof}

\begin{proposition}[Recurrence from one-cycle estimates]\label{prop:cycle-forest-recurrence}
Suppose that for each sufficiently large integer $r$, every connected simple graph $G$ of rank $r$
and every cycle $W$ with at most $\sqrt r$ edges satisfy
\[
 p(G,W)\le1/2+d(r),
\]
where $d(r)\ge0$.
Then, for all sufficiently large $r$,
\begin{equation}\label{eq:cycle-deletion-recurrence}
 \psi(r)\le\psi(r-1)+(1/4+d(r)/2)2^r.
\end{equation}
Consequently,
\begin{equation}\label{eq:cycle-recursion-window}
 \frac{\psi(r)}{2^r}
 \le\frac12+
 \max_{\lceil r/2\rceil\le s\le r}d(s)+2^{-r/2}.
\end{equation}
In particular, $d(r)=o(1)$ implies $\psi(r)\le(1/2+o(1))2^r$.
\end{proposition}

\begin{proof}
Choose a connected graph $G$ of rank $r$ with $c(G)=\psi(r)$.
Delete leaves, and suppress degree-two vertices.
For $r\ge2$, the resulting connected multigraph $K$ has minimum degree at least three.
These operations preserve its rank and cycle count, and $|V(K)|\le2r-2$.

The elementary Moore bound gives a cycle $W$ of length at most $2\log(2r)+1$ in $K$.
Loops and parallel pairs already give cycles of length one or two.
Otherwise, the usual vertex-centered and edge-centered breadth-first trees give the claimed logarithmic bound.
Choose an edge $e$ of $W$.
Subdivide every edge of $K$ twice, obtaining a simple graph $K'$ of the same rank.
The cycle $W$ becomes a cycle $W'$ with
\[
 |E(W')|\le6\log(2r)+3\le\sqrt r
\]
for sufficiently large $r$.
Subdivision preserves outside-forest probability and cycle counts.
Lemma~\ref{lem:recursion-marked-edge} therefore yields
\[
 c_e(K)\le(1/2+d(r))2^{r-1}.
\]
Since $e$ lies on a cycle, $K-e$ is connected and has rank $r-1$.
The deletion identity gives
\[
 \psi(r)=c(K-e)+c_e(K)
 \le\psi(r-1)+(1/2+d(r))2^{r-1},
\]
which proves \eqref{eq:cycle-deletion-recurrence}.

Put $z_r=\psi(r)/2^r-1/2$.
The recurrence becomes $z_r\le z_{r-1}/2+d(r)/2$.
For $s=\lfloor r/2\rfloor$, iteration gives
\[
 z_r\le2^{-(r-s)}z_s+
 \sum_{j=s+1}^{r}2^{j-r-1}d(j).
\]
Use $z_s\le1/2$ and sum the geometric weights, whose sum is less than one.
This proves \eqref{eq:cycle-recursion-window}.
\end{proof}

The manuscript's forest error $20/\log\log\log r$ gives the same order of cycle-count error through this recurrence.
The constant $21$ also follows for sufficiently large $r$.
If a one-cycle forest estimate has error $O(r^{-a})$ for a fixed $a>0$, the cycle-count error has the same order.
No additional term from cycles avoiding a selected witness set is needed.

\subsection{The additive excess as an edge problem}

For $r\ge2$, define
\[
 X_r=\psi(r)-2^{r-1},\qquad
 g(r)=\max_{\beta(G)=r}\ \min_{e\text{ on a cycle}}c_e(G)-2^{r-2}.
\]
The maximum may be taken over connected simple graphs.
It exists because the minimum edge frequency is an integer between one and $2^{r-1}$.

\begin{proposition}\label{prop:cycle-edge-excess-recurrence}
For every $r\ge2$,
\[
 X_r-X_{r-1}\le g(r).
\]
Hence, for $1\le s<r$,
\[
 X_r\le X_s+\sum_{j=s+1}^r g(j).
\]
\end{proposition}

\begin{proof}
Choose a graph attaining $\psi(r)$ and an edge $e$ of minimum frequency among its cyclic edges.
Then
\[
 \psi(r)=c(G-e)+c_e(G)
 \le\psi(r-1)+2^{r-2}+g(r).
\]
Subtract $2^{r-1}$ and iterate.
\end{proof}

For a simple graph and an edge $e=uv$, one has $c_e(G)=P_{uv}(G-e)$.
Thus an upper bound on $g(r)$ is a terminal-path bound for one suitable edge of each graph.
If $g(r)=O(2^r/r^a)$, then $X_r=O(2^r/r^a)$.
If $g(r)=O(2^{(1-\kappa)r})$ for $0<\kappa<1$, then $X_r$ has that same exponential order.
Both conclusions follow by summing the recurrence, splitting the sum at $r/2$ when needed.

The general additive stability theorem gives a fixed-gap bound, with a threshold depending on that gap.
It does not presently control this threshold strongly enough to give a polynomial error in arbitrary graphs.
Accordingly, the recurrence does not by itself prove such an error bound.

\subsection{An alternative proof of the fixed-minor error bound}

The recurrence also gives a proof of Theorem~\ref{thm:minor-cycle-error} without adding witness triangles.
We record it because it separates the geometric input from the counting step.

Fix $h\ge4$, put $K=4096h^3$, and set
\[
 A_h=8K(3K-1)2^{3K-2}.
\]
Let $\psi_h(r)$ be the maximum cycle count among simple $K_h$-minor-free graphs of rank $r$.
Suppose an extremizer has more than $2^{r-1}$ cycles.
The proof of Theorem~\ref{thm:minor-cycle-error} reduces it to a unique cyclic kernel $\Gamma$ with
\[
 3\le\deg(v)\le5,\qquad |V(\Gamma)|\le2r-2.
\]
Its underlying simple graph remains $K_h$-minor-free.

Choose a shortest cycle $W$ in $\Gamma$.
For $r\ge512$, the Moore bound gives
\[
 |V(W)|\le|E(W)|\le2\log(2r)+1\le r/16.
\]
The graph $J=\Gamma-V(W)$ therefore has rank at least $r-4|V(W)|$.
Apply Lemma~\ref{lem:path-small-components}, with the same $K$ as above, to delete at most $|V(J)|/16$ vertices.
Every remaining component has at most $K$ vertices, while the remaining rank is at least
\[
 r-4|V(W)|-4|V(J)|/16\ge r/2-4|V(W)|\ge r/4.
\]
Each component has rank at most $2K$.
Selecting one cycle from every cyclic component gives at least $r/(8K)$ disjoint cyclic regions outside $W$.
Their boundaries in $\Gamma$ are at most $3K$.
The sharper small-cut estimate \eqref{eq:path-small-cuts} gives
\[
 p(\Gamma,W)\le1/2+A_h/r.
\]

For any edge $e$ of $W$, Lemma~\ref{lem:recursion-marked-edge} now yields
\[
 c_e(\Gamma)\le(1/2+A_h/r)2^{r-1}.
\]
Choose an edge of the corresponding subdivided path in the original simple graph.
Deleting that edge removes exactly the cycles counted by $c_e(\Gamma)$ and decreases rank by one.
It preserves the excluded-minor condition.
Hence, whenever $\psi_h(r)>2^{r-1}$ and $r\ge512$,
\[
 \psi_h(r)\le\psi_h(r-1)+(1/2+A_h/r)2^{r-1}.
\]

For completeness, this recurrence gives the same explicit constant as Theorem~\ref{thm:minor-cycle-error}.
Put $C=A_h+1$.
For $r\le2C$, the trivial bound $\psi_h(r)<2^r$ gives
$\psi_h(r)\le(1/2+C/r)2^r$.
For larger $r$, assume the same bound at rank $r-1$.
If $\psi_h(r)\le2^{r-1}$, there is nothing to prove.
Otherwise the recurrence applies, and its normalized excess is at most
\[
 \frac{C}{2(r-1)}+\frac{A_h}{2r}\le\frac Cr.
\]
The last inequality is equivalent to $C/(r-1)\le1$, which holds in this range.
This proves the error bound by induction, without any witness addition.
